\documentclass[UTF8,fontset=fandol,12pt,a4paper]{ctexart}
\usepackage{amssymb, amsmath}
\usepackage{pgfplots}
\pgfplotsset{compat=1.18}

\title{信号与系统分析-hw3}
\author{黄泱博\quad2024012867}
\date{\today}

\begin{document}
\maketitle

\section*{3.8 确定满足给定条件的两个不同信号}

\noindent $x(t)$是实奇函数，则$a_k^*=a_{-k}=-a_k$，即有$Re(a_k)=0$ \\
由Parseval定理，$1=\frac{1}{2}\int_{0}^{2}|x(t)|^2dt=\sum_{k=-\infty}^{\infty}|a_k|^2=|a_{-1}|^2+|a_0|^2+|a_1|^2$\\
于是可以构造以下两个信号：
\[a_{-1}=-\frac{\sqrt2}{2}j,a_0=0,a_1=\frac{\sqrt2}{2}j, x_1(t)=-\frac{\sqrt2}{2}je^{-j\pi t}+\frac{\sqrt2}{2}je^{j\pi t}\]
\[a_{-1}=\frac{\sqrt2}{2}j,a_0=0,a_1=-\frac{\sqrt2}{2}j, x_2(t)=\frac{\sqrt2}{2}je^{-j\pi t}-\frac{\sqrt2}{2}je^{j\pi t}\]

\section*{3.14 已知输入输出求$H(e^{jk\pi /2})$在部分$k$处的取值}

$$x[n]=\sum_{k=-\infty}^{\infty}\delta[n-4k]$$

$$a_k=\frac{1}{N}\sum_{<N>}x[n]e^{-jk\frac{\pi}{2}n}=\frac{1}{4},\forall k$$

$$y[n]=cos(\frac{5\pi}{2}n+\frac{\pi}{4})=\frac{e^{j(\frac{5\pi}{2}n+\frac{\pi}{4})}+e^{-j(\frac{5\pi}{2}n+\frac{\pi}{4})}}{2}=\frac{e^{j\frac{\pi}{4}}}{2}e^{j\frac{5\pi}{2}n}+\frac{e^{-j\frac{\pi}{4}}}{2}e^{-j\frac{5\pi}{2}n}$$

\[b_3=b_{-5}=\frac{e^{-j\frac{\pi}{4}}}{2},b_1=b_5=\frac{e^{j\frac{\pi}{4}}}{2}\]

于是有：
$$b_k=a_k H(e^{jk\frac{\pi}{2}}) \Rightarrow H(e^{jk\frac{\pi}{2}})=4b_k$$

从而：
\[H(e^{j0\frac{\pi}{2}})=0,H(e^{j1\frac{\pi}{2}})=2e^{j\frac{\pi}{4}},H(e^{j2\frac{\pi}{2}})=0,H(e^{j3\frac{\pi}{2}})=2e^{-j\frac{\pi}{4}}\]

\section*{3.23 已知周期为4的Fourier级数，求信号}

\subsection*{(a)}

\[
a_k=\begin{cases}
  0, &k=0 \\
  {(j)}^k \frac{\sin{k\pi/4}}{k\pi} &else
\end{cases}
\]

\noindent $j^k=(e^{j\frac{\pi}{2}})^k=e^{jk\frac{\pi}{2}}$ 即为向左时移1个单位\\
已知$2T_1,T$的周期方波的系数$b_k=\frac{\sin{k\frac{2\pi}{T}T_1}}{k\pi}$,因此$a_k$对应于一个宽度为1，高度为1的方波。
又由于$0=a_0=\frac{1}{T}\int_{T}x(t)dt$，故需要向下平移$\frac{1}{4}$个单位。\\

\noindent 因此，在一个周期$(-3,1)$内，$x(t)=\begin{cases}
  \frac{3}{4} &-\frac{3}{2}<t<-\frac{1}{2} \\
  -\frac{1}{4} &-3<t<-\frac{3}{2}, -\frac{1}{2}<t<1
\end{cases}$

\subsection*{(b)}

\[a_k={(-1)}^k\frac{\sin{k\pi/8}}{2k\pi},a_0=\frac{1}{16}\]

\noindent 同上，${(-1)}^k=e^{jk\pi}$为向左时移2个单位，对应于一个宽度为$\frac{1}{2}$，高度为$\frac{1}{2}$的方波。又由于$\frac{1}{16}=a_0=\frac{1}{T}\int_{T}x(t)dt$,故不需要向下平移。

\noindent 因此，在一个周期$(0,4)$内，$x(t)=\begin{cases}
  \frac{1}{2} &\frac{7}{4}<t<\frac{9}{4} \\
  0, &0<t<\frac{7}{4}, \frac{9}{4}<t<4
\end{cases}$

\subsection*{(c)}

\[a_k=\begin{cases}
  jk &|k|<3\\
  0 &|k|\ge 3
\end{cases}\]

$$x(t)=\sum_{k=-\infty}^{\infty}a_ke^{jk\frac{\pi}{2}t}=\sum_{k=1}^{2}jk(e^{jk\frac{\pi}{2}t}-e^{-jk\frac{\pi}{2}t})=jk\times 2j\sin{k\frac{\pi}{2}t}=-2\sin{\frac{\pi t}{2}}-3sin{\pi t}$$

\subsection*{(d)}
$$a_k=\begin{cases}
  1 &k\equiv 0\mod 2 \\
  2 &k\equiv 1\mod 2
\end{cases}=\frac{3}{2}-\frac{1}{2}{(-1)}^k$$

\noindent 已知$\sum_{k=-\infty}^{\infty}\delta(t-kT)$的系数为$b_k=\frac{1}{T}$，同上有${(-1)}^k$为向左时移2个单位。
于是有$x(t)=6\sum_{k=-\infty}^{\infty}\delta(t-4k)-2\sum_{k=-\infty}^{\infty}\delta(t-4k-2)$

\section*{3.25}

$$x(t)=\cos{(4\pi t)}$$
$$y(t)=\sin{(4\pi t)}$$
$$z(t)=x(t)y(t)$$

\subsection*{(a) 求$x(t)$ Fourier级数}

$$x(t)=\cos{(4\pi t)}=\frac{e^{j4\pi t}+e^{-j4\pi t}}{2}$$
因此有$a_{-1}=a_1=\frac{1}{2}$，其余均为0

\subsection*{(b) 求$y(t)$ Fourier级数}

$$y(t)=\sin{(4\pi t)}=\frac{e^{j4\pi t}-e^{-j4\pi t}}{2j}$$
因此有$b_{-1}=-\frac{1}{2j},b_1=\frac{1}{2j}$,其余均为0

\subsection*{(c) 利用(a)(b),求$z(t)$ Fourier级数}

$$c_k=a_k*b_k=\sum_{l=-\infty}^{\infty}a_lb_{k-l}=a_{-1}b_{k+1}+a_1b_{k-1}$$
因此有$c_{-2}=-\frac{1}{4j},c_0=0,c_2=\frac{1}{4j}$,其余均为0

\subsection*{(d) 直接求$z(t)$ Fourier级数 并与(c)中结果比较}

$$z(t)=x(t)y(t)=\cos{(4\pi t)}\sin{(4\pi t)}=\frac{1}{2}sin{(8\pi t)}=\frac{1}{2}\frac{e^{j2\times 4\pi t}-e^{-j2\times 4\pi t}}{2j}$$
因此有$c_{-2}=-\frac{1}{4j},c_2=\frac{1}{4j}$,其余均为0，与(c)中结果一致

\section*{3.28 求 Fourier级数 并绘制$a_k$的模和相位}

\subsection*{(a) P3.28(a)-P3.28(c)}

\subsubsection*{(i)}
\[
N=7,\qquad \omega_0=\frac{2\pi}{N}=\frac{2\pi}{7}.
\]
\begin{align*}
  a_k
    &=\frac{1}{7}\sum_{n=0}^{6}x[n]e^{-jk\frac{2\pi}{7}n}
      =\frac{1}{7}\sum_{n=0}^{4}e^{-jk\frac{2\pi}{7}n} \\
    &=\begin{cases}
      \dfrac{5}{7}, & k\equiv0\pmod 7,\\[0.6em]
      \dfrac{1-e^{-j10\pi k/7}}{7\bigl(1-e^{-j2\pi k/7}\bigr)},
        & k\not\equiv0\pmod 7.
    \end{cases}
\end{align*}
\[
\boxed{
  a_k=\begin{cases}
    \dfrac{5}{7}, & k\equiv0\pmod 7,\\[0.6em]
    \dfrac{1}{7}e^{-j4\pi k/7}
      \dfrac{\sin(5\pi k/7)}{\sin(\pi k/7)},
      & k\not\equiv0\pmod 7.
  \end{cases}
}
\]

\begin{center}
\begin{minipage}[t]{0.48\textwidth}
\centering
\begin{tikzpicture}
  \begin{axis}[
    width=0.97\linewidth,height=5.3cm,
    title={模谱},
    axis lines=left,
    xlabel={$k$},ylabel={$|a_k|$},
    xmin=-0.5,xmax=6.6,ymin=0,ymax=0.82,
    xtick={0,1,...,6},ytick={0,0.2,0.4,0.6,0.8},
    tick label style={font=\small},
  ]
    \addplot[ycomb,blue!70!black,thick,mark=*,mark size=2pt]
      coordinates {
        (0,0.714285714286) (1,0.257419676544) (2,0.178139943388)
        (3,0.063577409702) (4,0.063577409702)
        (5,0.178139943388) (6,0.257419676544)
      };
  \end{axis}
\end{tikzpicture}
\end{minipage}\hfill
\begin{minipage}[t]{0.48\textwidth}
\centering
\begin{tikzpicture}
  \begin{axis}[
    width=0.97\linewidth,height=5.3cm,
    title={相位谱},
    axis lines=left,
    xlabel={$k$},ylabel={$\arg a_k$},
    xmin=-0.5,xmax=6.6,ymin=-2.2,ymax=2.2,
    xtick={0,1,...,6},
    ytick={-1.795195802,-0.897597901,0,0.897597901,1.795195802},
    yticklabels={$-\frac{4\pi}{7}$,$-\frac{2\pi}{7}$,$0$,
      $\frac{2\pi}{7}$,$\frac{4\pi}{7}$},
    tick label style={font=\small},
  ]
    \addplot[gray,thin] coordinates {(-0.5,0) (6.6,0)};
    \addplot[ycomb,blue!70!black,thick,mark=*,mark size=2pt]
      coordinates {
        (0,0) (1,-1.795195802051) (2,-0.448798950513)
        (3,0.897597901026) (4,-0.897597901026)
        (5,0.448798950513) (6,1.795195802051)
      };
  \end{axis}
\end{tikzpicture}
\end{minipage}
\end{center}

\subsubsection*{(ii)}

\[
N=6,\qquad\omega_0=\frac{2\pi}{N}=\frac{\pi}{3}.
\]
\begin{align*}
  a_k&=\frac{1}{6}\sum_{n=0}^{5}x[n]e^{-jk\frac{\pi}{3}n}=\frac{1}{6}\sum_{n=0}^{3}e^{-jk\frac{\pi}{3}n} \\
     &=\begin{cases}
      \dfrac{2}{3}, &k\equiv 0\pmod6,\\[0.6em]
      \dfrac{1-e^{-j4k\pi/3}}{6\bigl(1-e^{-jk\pi/3}\bigr)},
        &k\not\equiv0\pmod6.
     \end{cases}
\end{align*}
\[
\boxed{
  a_k=\begin{cases}
    \dfrac{2}{3}, &k\equiv0\pmod6,\\[0.6em]
    \dfrac{1}{6}e^{-jk\pi/2}
      \dfrac{\sin(2k\pi/3)}{\sin(k\pi/6)},
      &k\not\equiv0\pmod6.
  \end{cases}
}
\]
一个周期内的系数为
\[
  (a_0,a_1,a_2,a_3,a_4,a_5)
  =\left(\frac23,-\frac{\sqrt3}{6}j,\frac16,0,
    \frac16,\frac{\sqrt3}{6}j\right)
\]

\begin{center}
\begin{minipage}[t]{0.48\textwidth}
\centering
\begin{tikzpicture}
  \begin{axis}[
    width=0.97\linewidth,height=5.3cm,
    title={模谱},
    axis lines=left,
    xlabel={$k$},ylabel={$|a_k|$},
    xmin=-0.5,xmax=5.6,ymin=-0.03,ymax=0.82,
    xtick={0,1,...,5},ytick={0,0.2,0.4,0.6,0.8},
    tick label style={font=\small},
  ]
    \addplot[gray,thin] coordinates {(-0.5,0) (5.6,0)};
    \addplot[ycomb,blue!70!black,thick,mark=*,mark size=2pt]
      coordinates {
        (0,0.666666666667) (1,0.288675134595) (2,0.166666666667)
        (3,0) (4,0.166666666667) (5,0.288675134595)
      };
  \end{axis}
\end{tikzpicture}
\end{minipage}\hfill
\begin{minipage}[t]{0.48\textwidth}
\centering
\begin{tikzpicture}
  \begin{axis}[
    width=0.97\linewidth,height=5.3cm,
    title={相位谱},
    axis lines=left,
    xlabel={$k$},ylabel={$\arg a_k$},
    xmin=-0.5,xmax=5.6,ymin=-2.1,ymax=2.1,
    xtick={0,1,...,5},
    ytick={-1.570796327,0,1.570796327},
    yticklabels={$-\frac{\pi}{2}$,$0$,$\frac{\pi}{2}$},
    tick label style={font=\small},
  ]
    \addplot[gray,thin] coordinates {(-0.5,0) (5.6,0)};
    \addplot[ycomb,blue!70!black,thick,mark=*,mark size=2pt]
      coordinates {
        (0,0) (1,-1.570796326795) (2,0)
        (4,0) (5,1.570796326795)
      };
  \end{axis}
\end{tikzpicture}
\end{minipage}
\end{center}

\subsubsection*{(iii)}

\[
N=6,\qquad\omega_0=\frac{2\pi}{6}=\frac{\pi}{3}.
\]
\begin{align*}
  a_k&=\frac{1}{6}\sum_{n=0}^{5}x[n]e^{-jk\frac{\pi}{3}n}\\
     &=\boxed{\frac{1}{6}\left[
       1+4\cos\left(\frac{k\pi}{3}\right)
       -2\cos\left(\frac{2k\pi}{3}\right)\right]}.
\end{align*}
计算有：
\[
  a_0=\frac{1}{2},\quad a_1=\frac{2}{3},\quad a_2=0,\quad
  a_3=-\frac{5}{6},\quad a_4=0,\quad a_5=\frac{2}{3}.
\]

\begin{center}
\begin{minipage}[t]{0.48\textwidth}
\centering
\begin{tikzpicture}
  \begin{axis}[
    width=0.97\linewidth,height=5.3cm,
    title={模谱},
    axis lines=left,
    xlabel={$k$},ylabel={$|a_k|$},
    xmin=-0.5,xmax=5.6,ymin=-0.04,ymax=1,
    xtick={0,1,...,5},ytick={0,0.2,0.4,0.6,0.8,1},
    tick label style={font=\small},
  ]
    \addplot[gray,thin] coordinates {(-0.5,0) (5.6,0)};
    \addplot[ycomb,blue!70!black,thick,mark=*,mark size=2pt]
      coordinates {
        (0,0.5) (1,0.666666666667) (2,0)
        (3,0.833333333333) (4,0) (5,0.666666666667)
      };
  \end{axis}
\end{tikzpicture}
\end{minipage}\hfill
\begin{minipage}[t]{0.48\textwidth}
\centering
\begin{tikzpicture}
  \begin{axis}[
    width=0.97\linewidth,height=5.3cm,
    title={相位谱},
    axis lines=left,
    xlabel={$k$},ylabel={$\arg a_k$},
    xmin=-0.5,xmax=5.6,ymin=-0.4,ymax=3.7,
    xtick={0,1,...,5},
    ytick={0,1.570796327,3.141592654},
    yticklabels={$0$,$\frac{\pi}{2}$,$\pi$},
    tick label style={font=\small},
  ]
    \addplot[gray,thin] coordinates {(-0.5,0) (5.6,0)};
    \addplot[ycomb,blue!70!black,thick,mark=*,mark size=2pt]
      coordinates {(0,0) (1,0) (3,3.141592653590) (5,0)};
  \end{axis}
\end{tikzpicture}
\end{minipage}
\end{center}

\subsection*{(b)}
$$x[n]=\sin{(2\pi n/3)}\cos{(\pi n/2)}$$

\noindent 正弦项周期为 $3$，余弦项周期为 $4$，基波周期为
\[
  N=12,\qquad\omega_0=\frac{2\pi}{N}=\frac{\pi}{6}.
\]
在一个周期 $0\le n\le11$ 内，非零样值为
\[
  x[2]=x[4]=\frac{\sqrt3}{2},\qquad
  x[8]=x[10]=-\frac{\sqrt3}{2}.
\]
因此
\begin{align*}
  a_k&=\frac{1}{12}\sum_{n=0}^{11}x[n]e^{-jk\frac{\pi}{6}n}\\
     &=\frac{\sqrt3}{24}\left(
       e^{-jk\pi/3}+e^{-j2k\pi/3}
       -e^{-j4k\pi/3}-e^{-j5k\pi/3}\right).
\end{align*}
计算有
\[
\boxed{
  a_k=\begin{cases}
    -\dfrac{j}{4}, &k\equiv1,7\pmod{12},\\[0.5em]
    \dfrac{j}{4}, &k\equiv5,11\pmod{12},\\[0.5em]
    0, &\text{其他}.
  \end{cases}
}
\]

\begin{center}
\begin{minipage}[t]{0.48\textwidth}
\centering
\begin{tikzpicture}
  \begin{axis}[
    width=0.97\linewidth,height=5.3cm,
    title={模谱},
    axis lines=left,
    xlabel={$k$},ylabel={$|a_k|$},
    xmin=-0.5,xmax=11.7,ymin=-0.015,ymax=0.32,
    xtick={0,1,3,5,7,9,11},
    ytick={0,0.125,0.25},
    yticklabels={$0$,$\frac18$,$\frac14$},
    tick label style={font=\small},
  ]
    \addplot[gray,thin] coordinates {(-0.5,0) (11.7,0)};
    \addplot[ycomb,blue!70!black,thick,mark=*,mark size=2pt]
      coordinates {
        (0,0) (1,0.25) (2,0) (3,0) (4,0) (5,0.25)
        (6,0) (7,0.25) (8,0) (9,0) (10,0) (11,0.25)
      };
  \end{axis}
\end{tikzpicture}
\end{minipage}\hfill
\begin{minipage}[t]{0.48\textwidth}
\centering
\begin{tikzpicture}
  \begin{axis}[
    width=0.97\linewidth,height=5.3cm,
    title={相位谱},
    axis lines=left,
    xlabel={$k$},ylabel={$\arg a_k$},
    xmin=-0.5,xmax=11.7,ymin=-2.1,ymax=2.1,
    xtick={0,1,3,5,7,9,11},
    ytick={-1.570796327,0,1.570796327},
    yticklabels={$-\frac{\pi}{2}$,$0$,$\frac{\pi}{2}$},
    tick label style={font=\small},
  ]
    \addplot[gray,thin] coordinates {(-0.5,0) (11.7,0)};
    \addplot[ycomb,blue!70!black,thick,mark=*,mark size=2pt]
      coordinates {
        (1,-1.570796326795) (5,1.570796326795)
        (7,-1.570796326795) (11,1.570796326795)
      };
  \end{axis}
\end{tikzpicture}
\end{minipage}
\end{center}

\subsection*{(c)}

$x[n]$ 周期为 $4$，且
$x[n]=1-\sin(\pi n/4)$，$0\le n\le3$。
\[
  x[0]=1,\qquad x[1]=x[3]=1-\frac{\sqrt2}{2},\qquad x[2]=0.
\]
\begin{align*}
  a_k&=\frac{1}{4}\sum_{n=0}^{3}x[n]e^{-jk\frac{\pi}{2}n}\\
     &=\frac{1}{4}\left[1+\left(1-\frac{\sqrt2}{2}\right)
       \left(e^{-jk\pi/2}+e^{-j3k\pi/2}\right)\right]
\end{align*}
计算有
\[
  a_0=\frac{3-\sqrt2}{4},\qquad
  a_1=a_3=\frac{1}{4},\qquad
  a_2=\frac{\sqrt2-1}{4}
\]

\begin{center}
\begin{minipage}[t]{0.48\textwidth}
\centering
\begin{tikzpicture}
  \begin{axis}[
    width=0.97\linewidth,height=5.3cm,
    title={模谱},
    axis lines=left,
    xlabel={$k$},ylabel={$|a_k|$},
    xmin=-0.4,xmax=3.5,ymin=0,ymax=0.48,
    xtick={0,1,2,3},ytick={0,0.1,0.2,0.3,0.4},
    tick label style={font=\small},
  ]
    \addplot[ycomb,blue!70!black,thick,mark=*,mark size=2pt]
      coordinates {
        (0,0.396446609407) (1,0.25) (2,0.103553390593) (3,0.25)
      };
  \end{axis}
\end{tikzpicture}
\end{minipage}\hfill
\begin{minipage}[t]{0.48\textwidth}
\centering
\begin{tikzpicture}
  \begin{axis}[
    width=0.97\linewidth,height=5.3cm,
    title={相位谱},
    axis lines=left,
    xlabel={$k$},ylabel={$\arg a_k$},
    xmin=-0.4,xmax=3.5,ymin=-1.9,ymax=1.9,
    xtick={0,1,2,3},
    ytick={-1.570796327,0,1.570796327},
    yticklabels={$-\frac{\pi}{2}$,$0$,$\frac{\pi}{2}$},
    tick label style={font=\small},
  ]
    \addplot[gray,thin] coordinates {(-0.4,0) (3.5,0)};
    \addplot[ycomb,blue!70!black,thick,mark=*,mark size=2pt]
      coordinates {(0,0) (1,0) (2,0) (3,0)};
  \end{axis}
\end{tikzpicture}
\end{minipage}
\end{center}

\subsection*{(d)}
\noindent $x[n]$ 周期为 $12$，且
$x[n]=1-\sin(\pi n/4)$，$0\le n\le11$，故
\[
  N=12,\qquad\omega_0=\frac{\pi}{6}.
\]
令 $\theta=k\pi/6$，记
\begin{align*}
  S_k&=\sum_{n=0}^{11}\sin\left(\frac{\pi n}{4}\right)e^{-j\theta n}\\
     &=\frac{1}{2j}\left[
       \sum_{n=0}^{11}\left(e^{j\pi/4-j\theta}\right)^n
       -\sum_{n=0}^{11}\left(e^{-j\pi/4-j\theta}\right)^n
       \right].
\end{align*}
记
\[
  q_+=e^{j\pi/4-j\theta},\qquad q_-=e^{-j\pi/4-j\theta}
\]
等比数列求和有
\begin{align*}
  S_k&=\frac{1}{j}\left(\frac{1}{1-q_+}-\frac{1}{1-q_-}\right)\\
     &=\frac{1}{j}\frac{q_+-q_-}{(1-q_+)(1-q_-)}\\
     &=\frac{2e^{-j\theta}\sin(\pi/4)}
       {1-2e^{-j\theta}\cos(\pi/4)+e^{-j2\theta}}\\
     &=\frac{\sin(\pi/4)}{\cos\theta-\cos(\pi/4)}\\
     &=\frac{1}{\sqrt2\cos(k\pi/6)-1}.
\end{align*}
因此
\begin{align*}
  a_k&=\frac{1}{12}\sum_{n=0}^{11}
    \left[1-\sin\left(\frac{\pi n}{4}\right)\right]e^{-jk\pi n/6}\\
     &=\frac{1}{12}\sum_{n=0}^{11}e^{-jk\pi n/6}
       +\frac{1}{12\left[1-\sqrt2\cos(k\pi/6)\right]}.
\end{align*}

\[
\begin{aligned}
  a_0&=\frac{11-\sqrt2}{12},
    &a_6&=\frac{\sqrt2-1}{12},\\[0.4em]
  a_1=a_{11}&=-\frac{2+\sqrt6}{12},
    &a_2=a_{10}&=\frac{2+\sqrt2}{12},\\[0.4em]
  a_3=a_9&=\frac{1}{12},
    &a_4=a_8&=\frac{2-\sqrt2}{12},\\[0.4em]
  a_5=a_7&=\frac{\sqrt6-2}{12}.
\end{aligned}
\]

\begin{center}
\begin{minipage}[t]{0.48\textwidth}
\centering
\begin{tikzpicture}
  \begin{axis}[
    width=0.97\linewidth,height=5.3cm,
    title={模谱},
    axis lines=left,
    xlabel={$k$},ylabel={$|a_k|$},
    xmin=-0.5,xmax=11.7,ymin=0,ymax=0.92,
    xtick={0,1,3,5,7,9,11},ytick={0,0.2,0.4,0.6,0.8},
    tick label style={font=\small},
  ]
    \addplot[ycomb,blue!70!black,thick,mark=*,mark size=2pt]
      coordinates {
        (0,0.798815536469) (1,0.370790811899)
        (2,0.284517796864) (3,0.083333333333)
        (4,0.048815536469) (5,0.037457478565)
        (6,0.034517796864) (7,0.037457478565)
        (8,0.048815536469) (9,0.083333333333)
        (10,0.284517796864) (11,0.370790811899)
      };
  \end{axis}
\end{tikzpicture}
\end{minipage}\hfill
\begin{minipage}[t]{0.48\textwidth}
\centering
\begin{tikzpicture}
  \begin{axis}[
    width=0.97\linewidth,height=5.3cm,
    title={相位谱},
    axis lines=left,
    xlabel={$k$},ylabel={$\arg a_k$},
    xmin=-0.5,xmax=11.7,ymin=-0.4,ymax=3.7,
    xtick={0,1,3,5,7,9,11},
    ytick={0,1.570796327,3.141592654},
    yticklabels={$0$,$\frac{\pi}{2}$,$\pi$},
    tick label style={font=\small},
  ]
    \addplot[gray,thin] coordinates {(-0.5,0) (11.7,0)};
    \addplot[ycomb,blue!70!black,thick,mark=*,mark size=2pt]
      coordinates {
        (0,0) (1,3.141592653590) (2,0) (3,0) (4,0) (5,0)
        (6,0) (7,0) (8,0) (9,0) (10,0) (11,3.141592653590)
      };
  \end{axis}
\end{tikzpicture}
\end{minipage}
\end{center}


\section*{3.34 已知单位冲激响应$h(t)=e^{-4|t|}$,在各输入下分别求响应Fourier级数}
\[
H(s)=\int_{-\infty}^{\infty}h(t)e^{-st}dt=\int_{-\infty}^{0}e^{(4-s)t}dt+\int_{0}^{\infty}e^{-(4+s)t}dt=\frac{8}{16-s^2}
\]
\subsection*{(a)$x(t)=\sum_{n=-\infty}^{\infty}\delta(t-n)$}
$$T=1, \omega_0=\frac{2\pi}{T}=2\pi$$

\noindent $x(t)$Fourier级数:$a_k=\frac{1}{T}=1$\\
$y(t)$Fourier级数:$b_k=H(jk\omega_0)a_k=\frac{8}{16+4\pi^2k^2}$

\subsection*{(b)$x(t)=\sum_{n=-\infty}^{\infty}{(-1)}^n\delta(t-n)$}
$$T=2,\omega_0=\frac{2\pi}{T}=\pi$$
$x(t)$Fourier级数:
\[
a_k=\frac{1}{T}\int_{-\frac{1}{2}}^{\frac{3}{2}}x(t)e^{-jk\omega_0t}dt=\frac{1}{2}(1-e^{-jk\pi})
\]
$y(t)$Fourier级数:
\[
b_k=H(jk\omega_0)a_k=\frac{4(1-e^{-jk\pi})}{16+k^2\pi^2}
\]

\subsection*{(c)$x(t)$为图示周期方波}

$$T=1,T_1=\frac{1}{4},\omega_0=\frac{2\pi}{T}=2\pi$$

\noindent $x(t)$Fourier级数:$a_k=\frac{\sin{(k\omega_0T_1)}}{k\pi}=\frac{\sin{(\frac{k\pi}{2})}}{k\pi}, a_0=\frac{1}{2}$ \\
$y(t)$Fourier级数:$b_k=H(jk\omega_0)a_k=\frac{8\sin{(\frac{k\pi}{2})}}{k\pi(16+4k^2\pi^2)},b_0=H(0)a_0=\frac{1}{4}$

\section*{3.36 求输出Fourier级数}
$$y[n]-\frac{1}{4}y[n-1]=x[n]$$

\subsection*{(a)$x[n]=\sin{(\frac{3\pi}{4}n)}$}

$$T=8,\omega_0=\frac{2\pi}{T}=\frac{\pi}{4}$$

两侧同做Fourier级数展开，有：
\[
\sum_{k=-\infty}^{\infty}(1-\frac{1}{4}e^{-jk\omega_0})a_ke^{jk\omega_0n}=\frac{e^{j3\omega_0n}-e^{-j3\omega_0n}}{2j}
\]

于是有：
\[
a_3=\frac{\frac{1}{2j}}{1-\frac{1}{4}e^{-j\frac{3\pi}{4}}}
\]
\[
a_{-3}=-\frac{\frac{1}{2j}}{1-\frac{1}{4}e^{j\frac{3\pi}{4}}}
\]
\indent 其余均为0(单个周期内)

\subsection*{(b)$x[n]=\cos{(\frac{\pi}{4}n)}+2\cos{(\frac{\pi}{2}n)}$}

$$T=8, \omega=\frac{2\pi}{T}=\frac{\pi}{4}$$

两侧同做Fourier级数展开：
\[
\sum_{k=-\infty}^{\infty}(1-\frac{1}{4}e^{-jk\omega_0})a_ke^{jk\omega_0n}=\frac{e^{j\omega_0n}+e^{-j\omega_0n}}{2}+2\frac{e^{j2\omega_0n}+e^{-j2\omega_0n}}{2}
\]

于是有：
\[
a_1=\frac{\frac{1}{2}}{1-\frac{1}{4}e^{-j\frac{\pi}{4}}}
\]
\[
a_{-1}=\frac{\frac{1}{2}}{1-\frac{1}{4}e^{j\frac{\pi}{4}}}
\]

\[
a_2=\frac{1}{1-\frac{1}{4}e^{-j\frac{\pi}{2}}}
\]
\[
a_{-2}=\frac{1}{1-\frac{1}{4}e^{j\frac{\pi}{2}}}
\]
\indent 其余均为0(单个周期内)

\section*{3.44}
$$T=6, \omega_0=\frac{2\pi}{T}=\frac{\pi}{3}$$
$$x(t) \in \mathbb{R} \Longrightarrow a_k^*=a_{-k}$$
$$a_k=0, k=0, k>2 \Longrightarrow a_k=a_{-k}^*=0, \forall k<-2$$ 
$$x(t)=-x(t-3)\Longrightarrow a_k=-a_ke^{-jk\pi}\Longrightarrow a_k=0, k=2,-2$$
$$a_1\in\mathbb{R}^+\Longrightarrow a_{-1}=a_k^*\in\mathbb{R}$$

由Parseval定理：
$$|a_{-1}|^2+|a_1|^2=\frac{1}{6}\int_{-3}^{3}|x(t)|^2dt=\frac{1}{2}$$
\indent 于是有：
$$a_{-1}=a_1=\frac{1}{2}$$
\indent 从而有:
$$x(t)=a_{-1}e^{-j\omega_0t}+a_1e^{j\omega_0t}=\cos{(\frac{\pi}{3}t)}$$

\section*{3.47 $x(t)=\cos{(2\pi t)}$视作周期为3的信号，求Fourier级数}

$$T=3,\omega_0=\frac{2\pi}{T}=\frac{2\pi}{3}$$

\[
x(t)=\frac{e^{j2\pi t}+e^{-j2\pi t}}{2}=\frac{e^{j3\omega_0 t}+e^{-j3\omega_0 t}}{2}
\Rightarrow a_{-3}=a_3=\frac{1}{2}
\]
\indent 其余均为0

\section*{3.50}

$$T=8,\omega_0=\frac{2\pi}{T}=\frac{\pi}{4}$$

\begin{align*}
  x[n]&=\sum_{k=0}^{7}a_ke^{jk\omega_0n}\\
      &=\sum_{k=0}^{3}a_k\left(e^{jk\omega_0n}-e^{j(k+4)\omega_0n}\right)\\
      &=\sum_{k=0}^{3}a_ke^{jk\omega_0n}\left(1+(-1)^{n+1}\right).
\end{align*}
因此有：$x[n]=0, \forall n\equiv0\pmod2$
再由条件：$x[2n+1]=(-1)^n$即可确定

\[
  (x[0],x[1],\ldots,x[7])=(0,1,0,-1,0,1,0,-1).
\]
\begin{center}
\begin{tikzpicture}
  \begin{axis}[
    width=0.85\textwidth,height=5.5cm,
    axis lines=left,
    xlabel={$n$},ylabel={$x[n]$},
    xmin=-0.6,xmax=7.8,ymin=-1.35,ymax=1.45,
    xtick={0,1,...,7},ytick={-1,0,1},
    tick label style={font=\small},
  ]
    \addplot[gray,thin] coordinates {(-0.6,0) (7.8,0)};
    \addplot[ycomb,blue!70!black,thick,mark=*,mark size=2pt]
      coordinates {
        (0,0) (1,1) (2,0) (3,-1)
        (4,0) (5,1) (6,0) (7,-1)
      };
  \end{axis}
\end{tikzpicture}
\end{center}

\end{document}
